\chapter{Method}
\label{sec:design}
\lhead{\emph{Chapter 3. Method}}  % page header

This chapter states the hypotheses (Section~\ref{sec:hypotheses}), describes the data (Section~\ref{sec:data}) and the
measures (Section~\ref{sec:measures}), and sets out the testing protocol (Section~\ref{sec:protocol}).
Section~\ref{sec:threats} notes two weaknesses of the design.

\section{Hypotheses}
\label{sec:hypotheses}

Table~\ref{tab:predictions} sets out what imitation and continuity predict for each research question.
Under imitation, researchers entering a topic face the most uncertainty, so they should flock to hot topics, adopt
rising tools early and write follow-on work that benchmarks or evaluates what others have proposed. If review rewards
following, papers in crowded topics, papers with branded titles (a short coined name and a colon) and papers with
keywords near their peak should be accepted more often. Under continuity, incumbents should consolidate a topic's
existing body of work, newcomers should bring work from elsewhere, new tools should spread through teams that already
have ICLR work, and acceptance should track how close a paper stays to its team's earlier work.

Some rival explanations make the same predictions. Visible status, cumulative advantage and expertise
\citep{merton1968matthew,ibrahim2026bias,li2017expertise} also predict that acceptance is associated with a team's
record. Network diffusion, in which less connected people adopt later \citep{coleman1957diffusion}, predicts the same as
continuity for new tools, namely that teams new to ICLR take them up later.

\begin{table}[tb]
\centering
\footnotesize
\setlength{\tabcolsep}{3pt}
\renewcommand{\arraystretch}{1.15}
\begin{tabularx}{\textwidth}{@{}L{2.1cm} L{2.2cm} Y L{3.9cm}@{}}
\toprule
Route (reference point) & Field (RQ1) & Entry (RQ2) & Gate (RQ3) \\
\midrule
\textbf{Imitation} (others, now) & Convergence on shared tools & Entrants flock to hot topics, adopt rising tools early
and write follow-on work that benchmarks or evaluates others' methods & Acceptance responds to crowding (the registered
test predicted a penalty); a premium for branded titles and for keywords near their peak \\
\addlinespace
\textbf{Continuity} (own record) & Convergence on shared tools & Entrants bring work from other topics; incumbents
consolidate; new tools spread through teams with earlier ICLR work & Acceptance tracks closeness to the team's own
record \\
\bottomrule
\end{tabularx}
\caption[What each route predicts for the field, at entry and at the gate]{What each route predicts at the three
levels. Lock-in and efficient choice also predict convergence on shared tools, so only entry and the gate can tell the
routes apart.}
\label{tab:predictions}
\end{table}

Not all of these predictions were registered in advance. Table~\ref{tab:registered} lists the hypotheses that were
registered before any outcome was examined. They cover imitation's predictions for entry and the crowding test at the
gate, plus tests of how a topic's contribution mix changes as it grows. The registered crowding test predicted a
penalty for crowded topics, whereas imitation, if review rewards it, predicts a premium; a result of either sign would
therefore show review responding to the crowd. Continuity's predictions were written after the tests had been run and
are used only to organise the results. Two further hypotheses on the use of fashionable keywords (H2d and H3c) were
dropped before any estimate, because their measure failed its validation check.

\begin{table}[tb]
\centering
\footnotesize
\setlength{\tabcolsep}{4pt}
\renewcommand{\arraystretch}{1.1}
\begin{tabularx}{\textwidth}{@{}l l Y l@{}}
\toprule
ID & RQ & Prediction & Sign \\
\midrule
H1a--c & 1 & As a topic accumulates papers, its share of method papers falls and its shares of measurement and empirical
papers rise & $-$, $+$, $+$ \\
H1d & 1 & Where more of a topic's papers built on FMs in the previous year, fewer papers propose a new architecture &
$-$ \\
H1e & 1 & \ldots{} and more papers are empirical or evaluation studies & $+$ \\
\addlinespace
H2a & 2 & Hot topics have a lower share of incumbent papers (entrants flock to them) & $-$ \\
H2b & 2 & Newcomer papers are less often method papers & $-$ \\
H2c & 2 & Newcomer papers are more often measurement papers & $+$ \\
\addlinespace
H3a & 3 & Papers in more crowded topics are accepted less often & $-$ \\
H3b & 3 & This penalty is larger for method papers & $-$ \\
\bottomrule
\end{tabularx}
\caption[Registered hypotheses]{Hypotheses registered before any paper in the full corpus was labelled. All models
include topic and year fixed effects (Section~\ref{sec:protocol}).}
\label{tab:registered}
\end{table}

\section{Data}
\label{sec:data}

The corpus was built from official OpenReview submission records, which were downloaded on documented dates. It
contains 55,938 ICLR submissions from 2017 to 2026, including rejected and withdrawn papers. Titles, abstracts,
author keywords, author lists and decisions all come from these records.

Before the analysis began, about one fifth of the submissions were selected at random and set aside as a held-out
set. The remaining submissions form the \emph{discovery set}: the 43,480 papers with a usable abstract (at least 30
words). The discovery set was used to develop and run the registered tests. The \emph{held-out set} contains 11,099
papers with an abstract; 1,309 submissions whose links to other submissions were unresolved were kept out of both sets.
The held-out papers were not analysed until after the discovery tests had been registered and run, and they were then
used only for tests declared in advance. Because the two sets share conference years, reviewers and authors, the
held-out tests guard against overfitting by the analyst but not against patterns specific to a year or to a pool of
reviewers. For this reason the held-out analysis is described as a pre-registered out-of-sample replication, not a
certified untouched holdout. Table~\ref{tab:years} gives the number of papers by year.

\begin{table}[tb]
\centering
\footnotesize
\begin{tabular}{@{}lrr@{}}
\toprule
Year & Discovery & Held-out \\
\midrule
2017 & 384 & 105 \\
2018 & 770 & 212 \\
2019 & 1,248 & 309 \\
2020 & 2,070 & 500 \\
2021 & 2,360 & 603 \\
2022 & 2,649 & 704 \\
2023 & 3,872 & 995 \\
2024 & 5,765 & 1,445 \\
2025 & 8,951 & 2,284 \\
2026 & 15,411 & 3,942 \\
\midrule
All & 43,480 & 11,099 \\
\bottomrule
\end{tabular}
\caption[Papers by year]{Papers with a usable abstract by year. 69.3\% of discovery papers are from 2024--2026.}
\label{tab:years}
\end{table}

Different questions use different subsets. The field analyses use all discovery papers. The entry analyses use the
42,325 papers from 2019--2026 that have at least one author identifier. The acceptance tests use the 32,190 discovery
papers from 2018--2026 that were accepted or rejected (40.1\% were accepted) and the 8,185 held-out papers from the same
years that were accepted or rejected. Papers from 2017 serve only as earlier history, and the 9,986 withdrawn discovery
papers are left out of the acceptance tests. The field and entry analyses use no acceptance decisions.

Existing collections of ICLR data were checked but not used as the main source. The Re$^2$ dataset
\citep{zhang2025re2} was considered first. Its coverage of the main ICLR conference ends in 2023 and omits 2019, so this
study collected its own data.

\section{Measures}
\label{sec:measures}

\subsection{Topics}
A \emph{topic map} is one clustering of all papers into 100 topics. Each paper's title and abstract are turned into a
vector with SPECTER2, a text encoder trained on scientific papers \citep{cohan2020specter,singh2023scirepeval}, and the
vectors are grouped by spherical $k$-means with $k = 100$. A \emph{topic} is one cluster. A topic map can be fitted on
papers from any set of years and then applied to all papers. Unless another map is named, topics come from the map
fitted on all years. Alternative maps (other values of $k$, a map that weights each year equally, and maps fitted on
2017--2020 only) are used as robustness checks.

\subsection{Contribution types}
Each paper is given one primary contribution type by an LLM (Claude Haiku) that follows a fixed codebook. The types are
method (including a new architecture), theory, empirical study, dataset or benchmark, evaluation, and other. A
\emph{method paper} proposes a method. A \emph{measurement paper} contributes a dataset, benchmark or evaluation. An
\emph{FM paper} uses or adapts a large pre-trained model, such as an LLM; this label is coded with the same codebook.
Over all discovery papers, 75.1\% are method papers and 7.6\% are measurement papers.

The labels were checked in two ways. First, the same LLM labelled a sample of 2,174 discovery papers a second time, in
an independent run that could not see the first labels. The two runs agreed at Cohen's $\kappa = 0.93$ (97.5\%
agreement for method papers and 99.3\% for measurement papers). Second, the author checked a random sample of 100 labels
by hand against the codebook, without help from anyone else. About 80 of the 100 labels were clearly correct. The other
20 or so were borderline cases in which the author could not decide between two types either, and the LLM's label was
kept. No label in the sample was judged clearly wrong, but the borderline cases are not counted as agreement.

These checks show that the LLM is consistent with itself. They do not show how often its labels would match those of
trained human coders. \citet{ziems2024can} found that LLMs labelling social science text without task-specific
training often reach only fair agreement with human coders. A second human coder would be needed to measure this for
the present labels, and Chapter~\ref{sec:conclusion} lists it as future work.

\subsection{Entry profiles}
Entry status is defined relative to a topic and counts ICLR submissions only, with any outcome. In a \emph{newcomer
paper}, no author has an earlier ICLR submission in the paper's topic, following the topic-relative definition of
\citet{biliotti2024breaking}; otherwise it is an \emph{incumbent paper}, and an author with such a submission is a
\emph{topic insider}. A newcomer paper comes from a \emph{first-time team} if no author has any earlier ICLR
submission, and from a \emph{returning topic newcomer} if some author has one in another topic. Authors are matched by
their OpenReview identifiers, never by name. In 2021 ICLR moved from e-mail-based to profile-based identifiers, and the
two are not linked, so some earlier histories are hidden and some incumbents appear as newcomers. This pulls the entry
contrasts towards zero.

\subsection{Two distances}
Both distances are measured in the same SPECTER2 vector space, but from different reference points.
\emph{Closeness to the team's own record}, \Fmax{}, follows \citet{bao2025persistence}. For each author with earlier
ICLR submissions since 2021 (an \emph{experienced author}), it is the cosine similarity between the paper and the mean
of that author's earlier submissions; \Fmax{} is the highest value among the paper's experienced authors. Probable
earlier versions of the same paper are excluded. \emph{Distance from earlier ICLR work} is one minus the mean cosine
similarity to the ten most similar submissions from earlier years, expressed as a percentile within topic and year.

\subsection{Predictors measured from the past}
Every predictor uses only information that was available before the paper's year. \emph{Prior FM adoption} is the FM
share of the topic's papers in the previous year. \emph{Crowding} is the natural log of one plus the number of the
topic's earlier papers. A \emph{hot topic} is one in the top fifth of growth before the paper's year.

Abstracts were downloaded after the decisions, so some accepted papers carry their final version. Links to code
repositories, which are usually added after acceptance, appear in 17.6\% of accepted but only 0.5\% of rejected
2023--2026 abstracts, and so leak the outcome \citep{kaufman2012leakage}. Removing them and recomputing the text vectors
leaves the discovery estimates of the distance-based acceptance tests unchanged. This check could not be repeated for
held-out papers.

\section{Testing protocol}
\label{sec:protocol}

The hypotheses, models and a rule for support were registered in dated documents, kept with the analysis code, before
the full corpus had been labelled. Unless stated otherwise, the models are linear probability models with topic and
year fixed effects, and standard errors are clustered by topic. For the registered tests, acceptance decisions were read
only by pre-committed analysis code, in a single logged run \citep{nosek2018preregistration,vandenakker2021preregistration}.

Each registered test was re-run under pre-specified variants: other topic maps and numbers of topics, the map that
weights years equally, and subsets of years and of papers. Tests were grouped into families by research question for
the Holm correction for multiple testing \citep{holm1979simple}. A registered hypothesis counts as supported if its
estimate has the predicted sign with a Holm-adjusted $p < 0.05$ in its family, and keeps that sign in at least 75\% of
the counted variants and in both hindsight variants (the map that weights years equally, and the 2017--2020 map applied
to 2021--2026 papers).

Analyses that were not registered at the start (post hoc analyses) were declared in writing before they were run. For
each question, a one-sided test on the held-out papers was declared before any held-out estimate for that question was
computed. A declared held-out test is supported if its one-sided $p$ is below 0.05. Holm-adjusted values across the six
secondary held-out tests are reported alongside, but do not set the label. Because the held-out tests are one-sided, a
supported held-out test confirms the sign of an association; its size should not be compared with the discovery
estimate.

Tests that are not supported are described in one of two ways. A test is \emph{uninformative} when its estimate lies
below its minimum detectable effect at 80\% power (MDE80): it cannot tell a small association from none. It is
\emph{inconclusive} otherwise, or when no MDE80 is available. The evidence labels \lab{replicated}, \lab{held-out
supported} and \lab{exploratory} are defined in Chapter~\ref{sec:intro}.

\section{Weaknesses of the design}
\label{sec:threats}

Two weaknesses should be kept in mind when reading the results. First, the standard errors may be too small. In a
simulation in which there was no true effect but papers were related across topics, the main entry tests gave a
significant result 12.6\% and 10.0\% of the time at a nominal 5\% level, and a wild cluster bootstrap
\citep{cameron2008bootstrap} did not correct this. Second, keyword popularity and title style had been explored against
acceptance on the discovery papers before the hypotheses were registered. The discovery estimate for titles is
therefore not a test, and the keyword-timing tests probe a dimension that had already been seen.
